\documentclass[12pt]{article}
\usepackage[utf8]{inputenc}
\usepackage[T1]{fontenc}
\usepackage{hyperref}
\title{Unserdeutsch Language Resources for K-12 Curriculum Development}
\author{Compiled by Claude Sonnet 5 (claude-sonnet-5) for Andr\'as Kornai}
\date{}
\begin{document}
\maketitle

\section{Introduction}
Unserdeutsch (Rabaul Creole German, ISO 639-3 uln) is a German-based creole of Papua New Guinea, now spoken by older adults only. It is unusually well documented for its size, mainly through a DFG-funded project (2015--2021) led by P\'eter Maitz and Werner K\"onig, with Craig A. Volker as a cooperating partner. No school materials in the language were found.

\section{Primary Educational Resources}
None found. The corpus and papers below could serve as source material for teaching materials.

\section{Linguistic and Reference Materials}
\begin{itemize}
\item G\"otze et al. (2017), Documenting Unserdeutsch (Rabaul Creole German): a workshop report (University of Bern repository record): \url{https://boris-portal.unibe.ch/entities/publication/15927e7c-c909-42b2-970a-cec0b18ed488}
\item Same report, alternative record: \url{https://boris-portal.unibe.ch/handle/20.500.12422/60693}
\item G\"otze et al. 2017 (PDF, University of Bern): \url{https://www.germanistik.unibe.ch/e40649/e40653/e714347/e714513/Goetze_et_al_2017_ger.pdf}
\item Documenting Unserdeutsch: reversing colonial amnesia (Bern record): \url{https://boris.unibe.ch/121401/}
\item Documenting Unserdeutsch: reversing colonial amnesia (James Cook University record): \url{https://researchonline.jcu.edu.au/52477/}
\item Bern record: \url{https://boris.unibe.ch/121389}
\item DFG project details: \url{https://gepris.dfg.de/gepris/projekt/275623802}
\end{itemize}
Access to some of the full texts is restricted to registered users, according to the search summary.

\section{Religious and Community Resources}
No scripture translations or community organisation pages were found in the searches.

\section{Digital and Audio Resources}
\begin{itemize}
\item Archive for Spoken German (AGD), Unserdeutsch corpus page: \url{https://agd.ids-mannheim.de/UNSD_extern.shtml}
\item EXMARaLDA announcement, Unserdeutsch corpus available in the DGD: \url{https://exmaralda.org/?p=6015}
\end{itemize}
Per the search summaries, DGD version 2.21 (January 2024) includes the first release of the corpus, with 55 audio recordings, 52 transcripts and metadata on 49 speakers. Craig Volker's 1979--80 recordings are reportedly in the AGD's OZ corpus; access terms were not confirmed.

\section{Curriculum Development Context}
There are very few speakers, all older, and no known teaching programme. The corpus makes heritage or documentation-based materials possible. Any curriculum would need the agreement of the speaker community, which is partly dispersed (Australia, per the project literature). The search results do not describe how such a curriculum would be organised.

\end{document}